\documentclass[11pt]{article}
\usepackage[margin=1in]{geometry}
\usepackage{booktabs}
\usepackage{hyperref}
\usepackage{listings}
\usepackage{xcolor}

\lstset{
  basicstyle=\ttfamily\small,
  backgroundcolor=\color{gray!10},
  frame=single,
  framerule=0pt,
  breaklines=true,
  columns=fullflexible
}

\title{Running Photopic Optimizations on AWS}
\author{Mika Nystr\"om}
\date{March 2026}

\begin{document}
\maketitle

\section{Motivation}

The photopic luminous efficacy optimizer uses COBYLA (a derivative-free
constrained optimization algorithm) with hierarchical subdivision from
2 to 129 spectral parameters.  Each run takes 10--30 minutes on an
Apple M3 laptop.  The runs are embarrassingly parallel: each CRI
constraint level is completely independent.  For large parameter
sweeps (varying CCT, CRI, R9 constraints, or exploring finer spectral
resolution at 257 or 513 parameters), cloud execution is attractive.

\section{What Needs to Run}

Each optimization requires:
\begin{itemize}
\item The \texttt{photopic} binary (Modula-3, statically linked against M3 libraries)
\item The \texttt{photopic.scm} Scheme script (loaded at runtime)
\item System library: \texttt{libgmp} (GNU Multiple Precision Arithmetic)
\end{itemize}

A typical invocation:
\begin{lstlisting}
./photopic -run 2700 90 -100 -scm -scmfile photopic.scm
\end{lstlisting}
This maximizes luminous efficacy for a 2700\,K spectrum subject to
CRI\,$\geq$\,90 with R9 unconstrained.  Output files (\texttt{.res}
and \texttt{.dat}) are written to the current directory.

\section{One-Time Setup}

\subsection{Build Environment}

The CM3 compiler and m3utils tree must be built on Linux.  AMD64\_LINUX
is the primary supported target.

\begin{enumerate}
\item Bootstrap CM3 from a release binary:
\begin{lstlisting}
wget <cm3-release-url>/cm3-boot-AMD64_LINUX.tar.gz
cmake -S cm3-boot -B build -G Ninja \
      -DCMAKE_INSTALL_PREFIX=$HOME/cm3-install
cmake --build build && cmake --install build
export PATH=$HOME/cm3-install/bin:$PATH
\end{lstlisting}

\item Build CM3 from source:
\begin{lstlisting}
cd cm3
scripts/concierge.py full-upgrade --backend c all
\end{lstlisting}

\item Build m3utils:
\begin{lstlisting}
cd intel-async/async-toolkit/m3utils/m3utils
export CM3=yes
gmake
\end{lstlisting}

\item Verify:
\begin{lstlisting}
./photopic/AMD64_LINUX/photopic -run 2700 90 -100 \
    -scm -scmfile photopic/src/photopic.scm
\end{lstlisting}
\end{enumerate}

\subsection{AMI or Container}

For repeated use, bake the build into an AMI or Docker image:
\begin{lstlisting}
FROM ubuntu:24.04
RUN apt-get update && apt-get install -y libgmp-dev
COPY photopic /opt/photopic/
COPY photopic.scm /opt/photopic/
WORKDIR /opt/photopic
ENTRYPOINT ["./photopic"]
\end{lstlisting}

Since the M3 binary is statically linked (via \texttt{build\_standalone()}
in \texttt{m3overrides}), the container only needs \texttt{libgmp} as an
external dependency.

\section{Instance Selection}

The optimizer is single-threaded and CPU-bound.  Each run benefits from
high single-thread performance but does not use multiple cores.  The
natural parallelism is across runs (one core per CRI level).

\begin{table}[h]
\centering
\begin{tabular}{@{}llrrr@{}}
\toprule
Instance & Architecture & vCPUs & \$/hr & \$/hr (spot) \\
\midrule
\texttt{c7g.medium}   & ARM64 (Graviton3) & 1  & 0.034 & $\sim$0.013 \\
\texttt{c7g.2xlarge}   & ARM64 (Graviton3) & 8  & 0.272 & $\sim$0.108 \\
\texttt{c7g.8xlarge}   & ARM64 (Graviton3) & 32 & 1.088 & $\sim$0.435 \\
\texttt{c7a.2xlarge}   & AMD64 (EPYC)      & 8  & 0.308 & $\sim$0.123 \\
\texttt{c7a.16xlarge}  & AMD64 (EPYC)      & 64 & 2.462 & $\sim$0.985 \\
\bottomrule
\end{tabular}
\caption{Compute-optimized EC2 instances (US East, March 2026 pricing).
  Spot prices are approximate and fluctuate.}
\label{tab:instances}
\end{table}

ARM64 (Graviton) instances offer better price/performance and would
use the same architecture as the development laptop (Apple M3), though
a separate Linux build is still required.

\section{Cost Estimates}

\begin{table}[h]
\centering
\begin{tabular}{@{}lrrr@{}}
\toprule
Workload & CPU-hours & On-demand & Spot \\
\midrule
8 CRI levels, 129 params     & 2.7  & \$0.10 & \$0.04 \\
8 CRI levels, 513 params     & $\sim$22  & \$0.75 & \$0.30 \\
50 CRI/R9 combinations       & $\sim$17  & \$0.58 & \$0.23 \\
Full sweep: 5 CCTs $\times$ 50 combos & $\sim$83 & \$2.85 & \$1.14 \\
\bottomrule
\end{tabular}
\caption{Estimated costs on \texttt{c7g} instances.  Assumes 20 minutes
  per run at 129 parameters, scaling quadratically with parameter count
  for COBYLA.}
\label{tab:costs}
\end{table}

Even aggressive parameter sweeps cost only a few dollars.  The
dominant cost is engineering time for the one-time build setup.

\section{Execution}

\subsection{Simple: Single Instance}

For moderate workloads (up to $\sim$50 runs), a single multi-core
instance with GNU Parallel is sufficient:

\begin{lstlisting}
# Generate all combinations
for cct in 2700 3000 3500 4000 5000; do
  for cri in 60 70 80 82 85 90 95 98; do
    echo "$cct $cri"
  done
done > jobs.txt

# Run in parallel (one per core)
parallel -j $(nproc) --colsep ' ' \
  './photopic -run {1} {2} -100 -scm -scmfile photopic.scm' \
  < jobs.txt
\end{lstlisting}

Results appear in the working directory.  Transfer them back with
\texttt{scp} or \texttt{aws s3 cp}.

\subsection{Scalable: AWS Batch}

For larger sweeps or routine use, AWS Batch automates instance
provisioning:

\begin{enumerate}
\item Push the Docker image to ECR
\item Define a Batch compute environment (spot instances, \texttt{c7g} family)
\item Define a job definition referencing the container
\item Submit jobs programmatically:
\begin{lstlisting}
for cct in 2700 3000 3500 4000 5000; do
  for cri in 60 70 80 82 85 90 95 98; do
    aws batch submit-job \
      --job-name "photopic-${cct}-${cri}" \
      --job-queue photopic-queue \
      --job-definition photopic-job \
      --container-overrides \
        "command=[\"-run\",\"$cct\",\"$cri\",\"-100\",\
                  \"-scm\",\"-scmfile\",\"photopic.scm\"]"
  done
done
\end{lstlisting}
\item Collect results from S3 (add an upload step to the container entrypoint)
\end{enumerate}

This is more setup but handles hundreds of runs with automatic scaling
and spot instance management.

\section{Recommendations}

\begin{enumerate}
\item \textbf{Start simple.}  Build on a \texttt{c7g.2xlarge} spot instance
  (\$0.11/hr).  Keep it as a development box for iterating.  Run the
  shell script directly.
\item \textbf{Bake an AMI} once the build works, so you can launch fresh
  instances without rebuilding.
\item \textbf{Move to Batch} only if you find yourself running hundreds of
  combinations regularly.
\item \textbf{Use spot instances} throughout---the workload is fault-tolerant
  (just re-run interrupted jobs) and spot pricing is 60--70\% cheaper.
\end{enumerate}

\end{document}
